\section{From Pixel Area to Particle Masses: The Complete OPH Spectrum Derivation}\label{from-pixel-area-to-particle-masses-the-complete-oph-spectrum-derivation}

\begin{quote}
\textbf{Integrated derivation within this manuscript}: this part applies the OPH framework developed in Part I.

\textbf{Abstract.} This part documents the quantitative OPH program in its current implementation. The calculation uses two external inputs: the pixel area \(P \equiv a_{\text{cell}}/\ell_P^2 = 1.63094\), calibrated from the gauge coupling sector via the screen\textquotesingle s entropy-matching condition (§5.4 of Part I of this manuscript), and the total screen capacity \(\log(\dim\mathcal{H}) \sim 10^{122}\), which enters only in the cosmological/neutrino branch (§14 of Part V of this manuscript). Every other quantity appearing in the derivation chain is either a mathematical constant, a consequence of the gauge group (itself selected from a refinement-stable edge-sector category by Tannaka-Krein/DR reconstruction plus the extended MAR admissibility package, §6.1--§6.2 of Part I of this manuscript), or a derived structural integer (\(N_c = 3\), \(N_g = 3\), \(\varepsilon = 1/6\), \(\delta = 2/9\)). Outputs are classified by epistemic status: quantities flowing from the calibrated gauge couplings are calibration-sector consistency checks; the Higgs/top critical-surface branch is the main independent quantitative output; charged-lepton, flavor-texture, neutrino, and hadron branches are weaker phenomenological or downstream numerical continuations. A complete audit of all constants and the full epistemic classification appear in §1A--§1B below.
\end{quote}

\begin{center}\rule{0.5\linewidth}{0.5pt}\end{center}

\subsection{1. Inputs and Contract}\label{1-inputs-and-contract}

\subsubsection{1.1 Physical Inputs}\label{11-physical-inputs}

The entire derivation pipeline uses exactly \textbf{two} physical inputs:

{\def\LTcaptype{none} % do not increment counter
\begin{longtable}[]{@{}>{\RaggedRight\arraybackslash}p{0.240\textwidth}>{\RaggedRight\arraybackslash}p{0.240\textwidth}>{\RaggedRight\arraybackslash}p{0.240\textwidth}>{\RaggedRight\arraybackslash}p{0.240\textwidth}@{}}
\toprule\noalign{}
Input & Symbol & Value & Origin \\
\midrule\noalign{}
\endhead
\bottomrule\noalign{}
\endlastfoot
Pixel area constant & \(P \equiv a_{\text{cell}}/\ell_P^2\) & 1.63094 & \textbf{Calibrated} from gauge couplings via entropy matching (§5.4 of Part I; see §1B) \\
Screen capacity & \(\log(\dim\mathcal{H}_{\text{tot}})\) & \(\sim 10^{122}\) & Inferred from observed cosmological constant \\
\end{longtable}
}

\subsubsection{1.2 What Counts as "Derived"}\label{12-what-counts-as-derived}

Everything else entering the computation falls into one of three categories:

\begin{enumerate}
\def\labelenumi{\arabic{enumi}.}
\item
  \textbf{Mathematical constants}: \(\pi\), \(e\), \(\zeta(3)\), Casimir eigenvalues, group dimensions, \(\beta\)-function coefficients , all determined by the gauge group structure (itself selected from the refinement-stable edge-sector category via Tannaka-Krein/DR reconstruction together with the extended MAR admissibility package).
\item
  \textbf{Derived quantities}: the unification scale \(M_U\), the unified coupling \(\alpha_U\), the electroweak VEV \(v\), all gauge couplings at \(M_Z\), the \(\mathbb{Z}_6\) defect parameter \(\varepsilon = 1/6\), the Koide phase \(\delta = 2/9\), and the Froggatt-Nielsen integer exponents.
\item
  \textbf{Numerical-method parameters}: loop order for RG running (1-loop SM for critical surface, 4-loop MSbar for \(\Lambda_{\overline{\text{MS}}}\)), lattice sizes and statistics for hadron ratios. These affect precision but not the prediction logic.
\end{enumerate}

\subsubsection{1.3 No-Leakage Guarantee}\label{13-no-leakage-guarantee}

The computation code enforces a runtime mutation test: after computing all outputs, the PDG reference table is scrambled and outputs are recomputed. Any change triggers an assertion failure. This verifies that no PDG values leak into the computation pipeline beyond the initial \(P\) calibration. Implemented in \texttt{oph\_predict\_compare.py::\_assert\_pdg\_not\_used()}.

\begin{center}\rule{0.5\linewidth}{0.5pt}\end{center}

\subsection{1A. Complete Audit of All Constants}\label{1a-complete-audit-of-all-constants}

\textbf{Purpose.} This section catalogs \emph{every} named constant, integer, or coefficient appearing anywhere in the prediction code or derivation chain, and explains its origin. The goal is to demonstrate that nothing has been "smuggled in" , every value is either (I) one of the two external inputs of the current implementation, (II) a derived structural quantity from the OPH axioms and stated premises, (III) a mathematical/group-theoretic constant uniquely fixed by the gauge group, or (IV) a standard QFT result that any textbook reproduces from the Lagrangian.

\subsubsection{1A.1 The Physical Input}\label{1a1-the-physical-input}

{\def\LTcaptype{none} % do not increment counter
\setlength{\LTleft}{0pt}
\setlength{\LTright}{0pt}
\setlength{\tabcolsep}{1.2pt}
\renewcommand{\arraystretch}{1.05}
\footnotesize
\begin{longtable}[]{@{}>{\RaggedRight\arraybackslash}p{0.13\textwidth}>{\RaggedRight\arraybackslash}p{0.09\textwidth}>{\RaggedRight\arraybackslash}p{0.18\textwidth}>{\RaggedRight\arraybackslash}p{0.33\textwidth}>{\RaggedRight\arraybackslash}p{0.15\textwidth}@{}}
\toprule\noalign{}
Constant & Value & Code variable & Origin & Status \\
\midrule\noalign{}
\endhead
\bottomrule\noalign{}
\endlastfoot
\(P \equiv a_{\text{cell}}/\ell_P^2\) & 1.63094 & \texttt{P\_DEFAULT} & The pixel area: area of one UV cell on the holographic screen in Planck units. Extracted from the entropy-matching condition \(P/4 = \bar{\ell}_{\text{SU(2)}}(t_2) + \bar{\ell}_{\text{SU(3)}}(t_3)\) (§5.4 of Part I of this manuscript). This relation is \emph{derived} from the axioms; the \emph{numerical value} is fixed by requiring consistency with the observed gauge couplings (just as Newton\textquotesingle s constant \(G\) is the one free parameter of general relativity). & \textbf{Single free parameter} for particle physics \\
\(\log(\dim\mathcal{H}_{\text{tot}})\) & \(\sim 10^{122}\) & \texttt{LOG\_DIM\_H\_DEFAULT} & Total screen capacity = de Sitter entropy \(S_{dS} = 3\pi/(G\Lambda)\). Enters only for neutrino masses and cosmological constant (§14 of Part V of this manuscript). Derived from the screen area law (§5 of Part I of this manuscript). & \textbf{Second input} (cosmology sector only) \\
\end{longtable}
}

\subsubsection{1A.2 Structural Integers Derived from Axioms}\label{1a2-structural-integers-derived-from-axioms}

{\def\LTcaptype{none} % do not increment counter
\setlength{\LTleft}{0pt}
\setlength{\LTright}{0pt}
\setlength{\tabcolsep}{1.2pt}
\renewcommand{\arraystretch}{1.05}
\footnotesize
\begin{longtable}[]{@{}>{\RaggedRight\arraybackslash}p{0.13\textwidth}>{\RaggedRight\arraybackslash}p{0.12\textwidth}>{\RaggedRight\arraybackslash}p{0.21\textwidth}>{\RaggedRight\arraybackslash}p{0.28\textwidth}>{\RaggedRight\arraybackslash}p{0.14\textwidth}@{}}
\toprule\noalign{}
Constant & Value & Code variable & Derivation & Part I of this manuscript ref \\
\midrule\noalign{}
\endhead
\bottomrule\noalign{}
\endlastfoot
\(N_c\) (colors) & 3 & \texttt{N\_c\_DEFAULT} & Witten\textquotesingle s global SU(2) anomaly requires \(N_c + 1\) SU(2) doublets per generation to be even, so \(N_c\) must be odd. \(N_c = 1\) is excluded (SU(1) trivial). The Selection Axiom MAR (third component of the complexity vector) picks the smallest nontrivial value. \textbf{Result}: \(N_c = 3\). & Theorem 6.14 (§6.9); Part II of this manuscript §6 \\
\(N_g\) (generations) & 3 & \texttt{N\_g\_DEFAULT} & Three independent constraints narrow the window: (1) CP violation requires \(N_g \geq 3\) (CKM phases \((N_g-1)(N_g-2)/2 > 0\)). (2) Asymptotic freedom of SU(2) with one Higgs doublet requires \(N_g(N_c+1) < 43/2\), giving \(N_g \leq 5\). (3) MAR (fourth component of the complexity vector) picks the smallest value in \(\{3, 4, 5\}\). \textbf{Result}: \(N_g = 3\). & Proposition 6.9 (§6.4); Part II of this manuscript §7 \\
Gauge group & \((SU(3)\times SU(2)\times U(1))/\mathbb{Z}_6\) & --- & Tannaka-Krein / Doplicher-Roberts reconstruction from a refinement-stable edge-sector category yields a compact gauge group \(G\) (Theorem 6.1, §6.1). The Selection Axiom MAR then selects the SM factors once the admissible class includes one connected abelian charge factor: the minimal coupled carrier \(\mathbb{C}^3 \otimes \mathbb{C}^2\) enforces the product structure, the commutant argument fixes the connected abelian factor to a single \(U(1)\), and the \(\mathbb{Z}_6\) quotient follows from hypercharge quantization (Proposition 6.6). & §6.1--6.2; Part II of this manuscript \\
\(\beta_{\text{EW}}\) & 4 & \texttt{beta\_ew(N\_c)} & Number of SU(2) doublets per generation = \(N_c\) quark doublets + 1 lepton doublet = \(N_c + 1 = 4\). This is a counting consequence of the Witten anomaly analysis that already fixed \(N_c = 3\). & §6.9, §6.19 \\
\(\varepsilon\) (Z₆ defect) & \(1/6\) & \texttt{defect\_\allowbreak epsilon\_\allowbreak Z6()} & The SM gauge group has a \(\mathbb{Z}_6\) center quotient. Each unit of \(\mathbb{Z}_6\) defect insertion removes \(\Delta S = \ln 6\) nats of entropy (MaxEnt weighting, Assumption B). The resulting Boltzmann suppression per defect is \(e^{-\ln 6} = 1/6\). \textbf{This is topologically fixed} by the quotient structure, not chosen. & §6.18, §6.21 \\
\(\delta\) (Koide phase) & \(2/9\) & computed inline & The holonomy phase on generation space: \(\delta = \beta_{\text{EW}} \cdot Y_Q / N_g = (N_c+1)/(2 N_c N_g) = 4/18 = 2/9\). Here \(Y_Q = 1/(2N_c)\) is the quark-doublet hypercharge (fixed by anomaly cancellation) and \(N_g = 3\) is the number of generations. All ingredients are previously derived. Experimental extraction: \(\delta_{\text{exp}} = 0.2222248 \pm 0.0000063\), matching \(2/9 = 0.2222\ldots\) within \(0.4\sigma\). & Proposition 13.3 (§13 of Part V of this manuscript) \\
Quark exponents & \(n_u = (6,3,0)\), \(n_d = (6,4,2)\) & \texttt{derive\_\allowbreak integer\_\allowbreak vectors()} & Minimal SU(3) hierarchy structure: \(n_{u} = (2N_c, N_c, 0)\) and \(n_d = (2N_c, N_c+1, N_c-1)\). These are the unique sequences that (a) span the range \([0, 2N_c]\) in \(N_g\) steps, (b) have the top quark unsuppressed (\(n_t = 0\)), and (c) match the observed up/down mass ordering under CKM rotation. With \(N_c = 3\): \((6,3,0)\) and \((6,4,2)\). & §6.21 \\
Lepton exponents & \(n_e = (7,4,3)\) & \texttt{derive\_\allowbreak lepton\_\allowbreak exponents()} & Derived algorithmically: the Koide roots \(r_k^2\) (from \(\delta = 2/9\)) must satisfy \(r_i^2/r_j^2 \approx \varepsilon^{n_i - n_j}\) for all pairs. The code scans integer triples near \((N_g, N_g+1, \ldots)\) and selects the unique triple minimizing the log-ratio residuals. \textbf{Result}: \((n_\tau, n_\mu, n_e) = (3, 4, 7)\). & §6.21, §13 of Part V of this manuscript \\
CKM angles & \(s_{12} = \varepsilon\), \(s_{23} = \varepsilon^2\), \(s_{13} = \varepsilon^3\) & inline & Powers of \(\varepsilon = 1/6\): the CKM mixing angles scale as \(\varepsilon^{\lvert n_i - n_j \rvert}\) where the exponents are the generation-gap between up-type and down-type defect charges. This is the standard Froggatt-Nielsen scaling. & §6.21 \\
\(Q\) (Koide ratio) & \(2/3\) & inline & The \(\mathbb{Z}_3\) mode balance on generation space: the circulant Hermitian matrix \(\Phi = a I + b P + b^* P^2\) on 3 generations satisfies \(Q = (1 + 2\lvert b/a \rvert^2)/3 = 2/3\) when \(\lvert b/a \rvert = 1/\sqrt{2}\), i.e., when the singlet and charged \(\mathbb{Z}_3\) modes have equal norm. This is the MaxEnt equilibrium. & Theorem 13.2 (§13 of Part V of this manuscript) \\
\end{longtable}
}

\subsubsection{1A.3 Standard QFT Constants (Fixed by the Gauge Group + Matter Content)}\label{1a3-standard-qft-constants-fixed-by-the-gauge-group--matter-content}

These constants are \emph{not} inputs , they are uniquely determined mathematical consequences of the gauge group SU(3) \(\times\) SU(2) \(\times\) U(1) with \(N_g\) generations of the standard fermion representations. Any QFT textbook (e.g., Peskin \& Schroeder) derives them from the Lagrangian.

{\def\LTcaptype{none} % do not increment counter
\setlength{\LTleft}{0pt}
\setlength{\LTright}{0pt}
\setlength{\tabcolsep}{1.2pt}
\renewcommand{\arraystretch}{1.05}
\footnotesize
\begin{longtable}[]{@{}>{\RaggedRight\arraybackslash}p{0.20\textwidth}>{\RaggedRight\arraybackslash}p{0.16\textwidth}>{\RaggedRight\arraybackslash}p{0.20\textwidth}>{\RaggedRight\arraybackslash}p{0.32\textwidth}@{}}
\toprule\noalign{}
Constant & Value & Code location & What determines it \\
\midrule\noalign{}
\endhead
\bottomrule\noalign{}
\endlastfoot
MSSM \(\beta\)-coefficients & \((b_1, b_2, b_3) = (33/5, 1, -3)\) & \texttt{B\_MSSM} & The edge-sector computation (§6.17 of Part I of this manuscript) derives \(\beta\)-function shifts \(\Delta b \approx (2.49, 4.17, 4.01)\) from the heat-kernel vacuum-polarization weighting and the \(\mathbb{Z}_6\) quotient. These match the shifts from SM to MSSM to better than 1\%. The coefficients themselves are standard 1-loop group theory: \(b_i = \sum_R T(R_i)\) summed over all matter multiplets. Using MSSM content is a \emph{consequence} of the edge-sector matching, not an assumption. In the synchronized D11 branch they govern the UV segment above the synchronization scale \(\mu_{\mathrm{sync}}\). \\
SM 1-loop \(\beta\)-coefficients & standard SM set & \texttt{B\_SM} & Standard 1-loop coefficients for the SM with \(N_g = 3\), namely \((41/10,\,-19/6,\,-7)\). In the synchronized D11 implementation they govern the low-scale transport from \(m_Z\) up to the synchronization scale \(\mu_{\mathrm{sync}}\approx 6.8\times 10^{11}\,\mathrm{GeV}\), not the entire interval up to \(M_U\). \\
4-loop MSbar \(\beta\)-coefficients & \(\beta_0, \beta_1, \beta_2, \beta_3\) & \texttt{beta\_\allowbreak coeffs\_\allowbreak msbar(n\_\allowbreak f)} in \texttt{oph\_\allowbreak qcd.py} & Standard perturbative QCD: \(\beta_0 = 11 - 2n_f/3\), \(\beta_1 = 102 - 38n_f/3\), etc. These are computed from SU(3) Feynman diagrams at each loop order. The 4-loop coefficient \(\beta_3\) involves \(\zeta(3)\); all are universal in the MSbar scheme. \\
Pole-mass conversion & \(K_2 = 13.44 - 1.04 n_l\), \(K_3 = 190.6 - 26.7 n_l + 0.65 n_l^2\) & \texttt{top\_\allowbreak pole\_\allowbreak from\_\allowbreak msbar()} & Standard 3-loop QCD relation between \(\overline{\text{MS}}\) and pole mass (Chetyrkin, Kniehl, Steinhauser 2000). These are perturbative coefficients computed from Feynman diagrams; they depend only on \(n_l\) (number of light flavors) and SU(3) group factors. \\
Casimir eigenvalues & \(C_2(p,q)\), \(d_{(p,q)}\) & \texttt{ellbar\_\allowbreak su3()}, \texttt{ellbar\_\allowbreak su2()} & Group-theoretic: \(C_2(p,q) = \frac{1}{3}(p^2 + q^2 + pq + 3p + 3q)\) for SU(3), \(C_2(j) = j(j+1)\) for SU(2), \(d_{(p,q)} = \frac{1}{2}(p+1)(q+1)(p+q+2)\). These are properties of the Lie algebra, not physical inputs. \\
\(\zeta(3)\) & 1.20206... & \texttt{z3} in \texttt{oph\_\allowbreak qcd.py} & Ap\'{e}ry\textquotesingle s constant , a pure mathematical constant appearing in 4-loop \(\beta_3\). \\
GUT normalization & \(\alpha_1 = \frac{5}{3}\alpha_Y\) & inline & The factor \(5/3\) is the standard GUT normalization ensuring all three couplings unify. It follows from embedding U(1)\(_Y\) into SU(5) and is fixed by the hypercharge assignments of the SM fermions. \\
\end{longtable}
}

\subsubsection{1A.4 Dimensional/Definitional Constants}\label{1a4-dimensionaldefinitional-constants}

{\def\LTcaptype{none} % do not increment counter
\setlength{\LTleft}{0pt}
\setlength{\LTright}{0pt}
\setlength{\tabcolsep}{1.2pt}
\renewcommand{\arraystretch}{1.05}
\footnotesize
\begin{longtable}[]{@{}>{\RaggedRight\arraybackslash}p{0.20\textwidth}>{\RaggedRight\arraybackslash}p{0.16\textwidth}>{\RaggedRight\arraybackslash}p{0.20\textwidth}>{\RaggedRight\arraybackslash}p{0.32\textwidth}@{}}
\toprule\noalign{}
Constant & Value & Code variable & Status \\
\midrule\noalign{}
\endhead
\bottomrule\noalign{}
\endlastfoot
\(E_P\) (Planck energy) & \(1.220890 \times 10^{19}\) GeV & \texttt{E\_PLANCK\_GEV} & \textbf{Definition}: \(E_P = \sqrt{\hbar c^5/G}\). This is a unit conversion factor, not a physical input , it converts from natural (Planck) units to GeV. In the OPH framework, \(G\) itself is derived from the pixel area via \(G = a_{\text{cell}}/(4\bar{\ell}\,\hbar/c^3)\) (§5.4 of Part I of this manuscript), so \(E_P\) is ultimately set by \(P\). \\
\(e^{2\pi}\) in \(M_U\) formula & 535.49... & inline & The Euclidean regularity condition on the collar geometry fixes the angular period to \(2\pi\) (§5.8 of Part I of this manuscript). This is a geometric constant, not a parameter. \\
\(\sqrt{2}\) in Yukawa & \(v/\sqrt{2}\) & inline & Standard Higgs mechanism convention: the Yukawa coupling \(y_f\) relates to the fermion mass via \(m_f = y_f v / \sqrt{2}\). This is a normalization convention, not physics. \\
\(\Delta\rho_{\text{stage-3}}\) & \(3y_t^2/(32\pi^2)\), evaluated at \(y_t \approx 1\) & inline & Universal 1-loop custodial correction with explicit Yukawa dependence. The displayed stage-3 number uses the unit-Yukawa estimate \(y_t \approx 1\). \\
\end{longtable}
}

\subsubsection{1A.5 Numerical-Method Parameters}\label{1a5-numerical-method-parameters}

These affect computational precision but not the prediction logic. Changing them changes the last digits, not the physics.

{\def\LTcaptype{none} % do not increment counter
\begin{longtable}[]{@{}>{\RaggedRight\arraybackslash}p{0.320\textwidth}>{\RaggedRight\arraybackslash}p{0.320\textwidth}>{\RaggedRight\arraybackslash}p{0.320\textwidth}@{}}
\toprule\noalign{}
Parameter & Value used & Effect of changing \\
\midrule\noalign{}
\endhead
\bottomrule\noalign{}
\endlastfoot
Loop order for RG (critical surface) & 1-loop UV-synchronized SM/MSSM piecewise transport & 2-loop running or a changed synchronization prescription shifts \(m_H\) and \(m_t\) at the GeV level \\
Loop order for \(\Lambda_{\overline{\text{MS}}}\) & 4-loop MSbar & 3-loop changes \(\Lambda\) by \textasciitilde2\% \\
RK4 step count & 2000 steps & Doubling changes nothing to displayed precision \\
Lattice sizes (hadrons) & \(L = 2\)--\(6\) & Larger volumes needed for precision; current values are prototype \\
Bisection tolerance & \(10^{-10}\) & Affects last digits of \(\alpha_U\) only \\
Heat-kernel truncation & 200 representations & Adding more changes \(\bar{\ell}\) by \(< 10^{-12}\) \\
\end{longtable}
}

\subsubsection{1A.6 Summary: What Is and Isn\textquotesingle t Assumed}\label{1a6-summary-what-is-and-isnt-assumed}

\textbf{Inherited from earlier Part I structural results (canonical package plus the stated gauge, heat-kernel, and refinement premises):}

\begin{itemize}
\tightlist
\item
  Gauge group structure (Tannaka-Krein reconstruction from edge sectors)
\item
  \(\varepsilon = 1/6\) (topological, from \(\mathbb{Z}_6\) quotient)
\item
  \(\lambda(M_U) = 0\), \(\beta_\lambda(M_U) = 0\) (refinement stability)
\item
  \(Q = 2/3\) (MaxEnt on \(\mathbb{Z}_3\) generation space)
\item
  Heat-kernel form \(p_R \propto d_R e^{-tC_2(R)}\) (MaxEnt + edge completion)
\end{itemize}

\textbf{Derived under the Selection Axiom MAR (Part II of this manuscript)}:

\begin{itemize}
\tightlist
\item
  SM gauge group \(SU(3) \times SU(2) \times U(1)/\mathbb{Z}_6\) (MAR + admissibility)
\item
  \(N_c = 3\) (Witten anomaly + MAR minimality)
\item
  \(N_g = 3\) (CP violation + asymptotic freedom + MAR minimality)
\item
  \(\delta = 2/9\) (algebraic consequence of \(N_c = 3\), \(N_g = 3\), \(\beta_{\text{EW}} = 4\))
\item
  \(\beta_{\text{EW}} = 4\) (doublet counting from \(N_c = 3\))
\item
  Integer exponents \(n_u\), \(n_d\), \(n_e\) (algorithmically from \(N_c\), \(N_g\), \(\varepsilon\), \(\delta\))
\end{itemize}

\textbf{Fixed by the gauge group (standard QFT, no free parameters)}:

\begin{itemize}
\tightlist
\item
  All \(\beta\)-function coefficients (1-loop through 4-loop)
\item
  Pole-mass conversion coefficients \(K_2\), \(K_3\)
\item
  Casimir eigenvalues, representation dimensions
\item
  GUT normalization factor \(5/3\)
\end{itemize}

\textbf{The primary particle-physics calibration input}:

\begin{itemize}
\tightlist
\item
  \(P = 1.63094\) (pixel area), calibrated from the gauge coupling sector (see §1B)
\end{itemize}

\textbf{Second external input (cosmology/capacity branch)}:

\begin{itemize}
\tightlist
\item
  \(\log(\dim\mathcal{H}) \sim 10^{122}\) (screen capacity)
\end{itemize}

\textbf{What is NOT assumed or smuggled in}:

\begin{itemize}
\tightlist
\item
  No particle masses
\item
  No CKM matrix elements
\item
  No Higgs VEV or quartic coupling
\item
  No SUSY-breaking scale
\item
  No Yukawa couplings
\item
  No \(\Lambda_{\text{QCD}}\)
\end{itemize}

\begin{center}\rule{0.5\linewidth}{0.5pt}\end{center}

\subsection{1B. Calibration vs.\ Prediction: Epistemic Classification of Outputs}\label{1b-calibration-vs-prediction}

This part is not the recovered core of OPH. In the master ledger of Part~I and the compact note, the recovered core is D1--D5 together with D7--D9. The present spectrum part primarily houses the current supplement-backed numerical implementation (D10--D11) and several D12 continuations. The purpose of this section is therefore not to inflate the quantitative program, but to state exactly which numbers are calibration-sector checks, which are secondary quantitative outputs, and which remain downstream or phenomenological branches.

\subsubsection{1B.1 The Calibration Structure}\label{1b1-the-calibration-structure}

The pixel area constant \(P\) is the external continuous input used by the present particle-physics implementation. In practice it is fixed through the gauge sector via the entropy-matching / pixel-constraint closure. Quantities algebraically entangled with that same closure sit in D10 and are calibration-sector consistency checks rather than independent confirmations. In this file, D10 includes the unified coupling and scale-setting branch, the running gauge couplings at \(m_Z\), \(\alpha_{\mathrm{em}}^{-1}(m_Z)\), \(\sin^2\theta_W(m_Z)\), \(v\), and the tree-level electroweak masses. Agreement there shows that the current implementation is self-consistent once \(P\) and the printed supplement-backed conventions are imposed. It does not by itself promote the whole spectrum program to theorem level.

\subsubsection{1B.2 Why ``Downstream of \(P\)'' Does Not Mean ``Recovered Core''}\label{1b2-dimensionless-independence}

Using \(P\) to set an overall energy scale and using an output branch as a recovered-core claim are different things. The critical-surface Higgs/top branch illustrates the distinction: once the D10 trajectory is fixed, the extra condition
\[
\lambda(M_U)=0,\qquad \beta_\lambda(M_U)=0
\]
adds no new continuous fit, so its failure would challenge the D11 branch rather than the D10 calibration closure. But it is still not part of the recovered core D1--D5 plus D7--D9. By contrast, charged-lepton/Koide fits, texture exponents, neutrino estimates, proton-spin bookkeeping, and downstream hadron extractions use additional ans\"atze, capacity assumptions, or external lattice/QCD machinery. Those belong to D12 or to later downstream numerical layers and must not be read back into the core claim set.

Put differently: scale-setting dependence does not erase all distinctions, but neither does it automatically elevate later numerics to theorem level. The correct sync rule is the master ledger: D10 is calibration-sector, D11 is the secondary quantitative branch, and the remaining spectrum branches are deferred continuations unless explicitly promoted elsewhere.

\subsubsection{1B.3 Classification}\label{1b3-classification}

{\footnotesize
\def\LTcaptype{none}
\begin{longtable}[]{@{}>{\RaggedRight\arraybackslash}p{0.22\textwidth}>{\RaggedRight\arraybackslash}p{0.28\textwidth}>{\RaggedRight\arraybackslash}p{0.40\textwidth}@{}}
\toprule
Tier / node & Representative outputs in this part & Status rule \\
\midrule
\endhead
\bottomrule
\endfoot
Phase I recovered core inherited from Parts I--II (D1--D5, D7--D9) & conditional relativity branch; compact gauge reconstruction in the bosonic branch; realized Standard Model quotient, exact hypercharges, \(N_c=3\), \(N_g=3\), and no gauge-mediated proton decay & These are structural or realized-branch outputs proved upstream. They are \emph{not} generated by the present spectrum numerics, even when this part reuses them. \\
Phase II input-dependent corollary inherited from Part~I (D6) & \(\Lambda=\frac{3\pi}{G N_{\mathrm{scr}}}\) once the screen-capacity identification is adopted, together with the separate order-of-magnitude neutrino estimate tied to the same capacity input & Separate from the recovered core. It uses the second external continuous input \(N_{\mathrm{scr}}\). \\
Phase II current implementation: calibration sector (D10) & pixel-closure gauge/electroweak consistency, \(\alpha_i(m_Z)\), \(\alpha_{\mathrm{em}}^{-1}(m_Z)\), \(\sin^2\theta_W(m_Z)\), \(v\), \(m_W\), and the running \(Z\)-mass relation & Supplement-backed consistency checks of the current implementation. Their agreement is expected once \(P\) and the printed running/matching conventions are imposed. \\
Phase II current implementation: secondary quantitative branch (D11) & Higgs/top critical-surface branch & Main non-calibration numerical branch of the present implementation. It adds no further continuous fit once D10 is fixed, but it remains secondary to the recovered core. \\
Phase III deferred continuations / downstream numerical branches (D12 and beyond) & charged-lepton/Koide fits, quark-texture branches, neutrino estimates beyond the D6 boundary, downstream hadron extractions, proton-spin bookkeeping, and any use of \(\varepsilon=1/6\) or \(\delta=2/9\) as flavor ans\"atze & Failure retracts the corresponding continuation only. These branches use extra structural choices, capacity assumptions, or external numerical machinery beyond D10--D11. \\
\end{longtable}
}

\subsubsection{1B.4 What This Part Actually Contributes}\label{1b4-what-oph-achieves}

The present spectrum part should therefore be read as follows. Its documentary center of gravity is the D10--D11 implementation layer. The gauge-sector closure and electroweak scale-setting branch are calibration-sector consistency checks. The Higgs/top critical-surface branch is the main secondary quantitative output. Charged-lepton, texture, neutrino, proton-spin, and hadronic branches remain explicitly weaker continuations.

This means that quantities sometimes reused in these later branches do \emph{not} all have the same claim tier. In particular, \(\varepsilon=1/6\), \(\delta=2/9\), Koide relations, texture exponents, neutrino scales, proton-spin fits, and hadron fits should not be cited as recovered-core outputs of overlap consistency plus the stated scaling-limit and categorical premises. The correct compact summary is simpler: Parts~I--II carry the recovered relativity-plus-Standard-Model core; this part begins at the D10 implementation layer and keeps everything beyond D11 explicitly non-core.

\begin{center}\rule{0.5\linewidth}{0.5pt}\end{center}

\subsection{2. Stage 1: Fundamental Scales}\label{2-stage-1-fundamental-scales}

\subsubsection{2.1 Unification Scale}\label{21-unification-scale}

The OPH framework derives a unification scale from the pixel area:

\[M_U = \frac{E_P}{e^{2\pi}} \cdot P^{1/6}\]

where \(E_P = 1.220890 \times 10^{19}\) GeV is the Planck energy. This gives:

\[M_U \approx 2.474 \times 10^{16} \text{ GeV}\]

\textbf{Derivation}: The factor \(e^{2\pi}\) arises from the modular geometry of the collar (the Euclidean regularity condition fixes the angular period to \(2\pi\)). The \(P^{1/6}\) factor comes from the dimensional relation between pixel area and the UV cell scale (§5.8 of Part I of this manuscript).

\subsubsection{2.2 Cell Energy Scale}\label{22-cell-energy-scale}

The UV cell energy is:

\[E_{\text{cell}} = \frac{E_P}{\sqrt{P}} \approx 9.56 \times 10^{18} \text{ GeV}\]

This is the natural energy scale of a single computational element on the holographic screen.

\begin{center}\rule{0.5\linewidth}{0.5pt}\end{center}

\subsection{3. Stage 2: Gauge Closure and the Pixel Constraint}\label{3-stage-2-gauge-closure-and-the-pixel-constraint}

\subsubsection{3.1 The Heat-Kernel Entropy}\label{31-the-heat-kernel-entropy}

The edge-center completion (Theorem 2.3 of Part I of this manuscript) combined with MaxEnt yields sector probabilities of the heat-kernel form:

\[p_R(t) \propto d_R \, e^{-t \, C_2(R)}\]

where \(d_R\) is the representation dimension, \(C_2(R)\) the quadratic Casimir, and \(t = 4\pi^2 \alpha\) is the diffusion parameter encoding the gauge coupling.

The \textbf{mean entropy per cell} (the \(\bar{\ell}\) function) for each gauge factor is:

\[\bar{\ell}_G(t) = \sum_R p_R(t) \log d_R\]

For SU(2):

\[\bar{\ell}_{\text{SU(2)}}(t) = \sum_{j=0,1/2,1,...} p_j(t) \log(2j+1), \quad p_j \propto (2j+1) e^{-t \, j(j+1)}\]

For SU(3):

\[\bar{\ell}_{\text{SU(3)}}(t) = \sum_{p,q \geq 0} p_{(p,q)}(t) \log d_{(p,q)}, \quad p_{(p,q)} \propto d_{(p,q)} \, e^{-t \, C_2(p,q)}\]

where \(d_{(p,q)} = \frac{1}{2}(p+1)(q+1)(p+q+2)\) and \(C_2(p,q) = \frac{1}{3}(p^2 + q^2 + pq + 3p + 3q)\).

\subsubsection{3.2 The Pixel Constraint}\label{32-the-pixel-constraint}

The generalized entropy matching (§5.4 of Part I of this manuscript) gives:

\[\frac{P}{4} = \bar{\ell}_{\text{SU(2)}}(t_2) + \bar{\ell}_{\text{SU(3)}}(t_3)\]

where \(t_i = 4\pi^2 \alpha_i\). This is the \textbf{pixel constraint}: the total edge entropy per cell must equal \(P/4\).

\subsubsection{\texorpdfstring{3.3 Gauge Running and the \(\alpha_U\) Solution}{3.3 Gauge Running and the \textbackslash alpha\_U Solution}}\label{33-gauge-running-and-the-alpha_u-solution}

At the unification scale, all three gauge couplings emerge from a single \(\alpha_U\). Running down to a scale \(\mu\) with MSSM-like \(\beta\)-function coefficients \((b_1, b_2, b_3) = (33/5, 1, -3)\):

\[\alpha_i^{-1}(\mu) = \alpha_U^{-1} + \frac{b_i}{2\pi} \ln\frac{M_U}{\mu}\]

\textbf{Why MSSM coefficients?} The edge-sector computation (§6.17 of Part I of this manuscript) derives \(\beta\)-function shifts \(\Delta b \approx (2.49, 4.17, 4.01)\) from the heat-kernel form, the \(\mathbb{Z}_6\) quotient structure, and the Peter-Weyl vacuum-polarization weighting. These match the MSSM shifts \((2.5, 4.17, 4.0)\) to better than 1\%. SM-only coefficients catastrophically fail (\(\alpha_s\) prediction off by 52\(\sigma\)).

\subsubsection{3.4 Self-Consistent Fixed Point}\label{34-self-consistent-fixed-point}

The system is solved self-consistently:

\begin{enumerate}
\def\labelenumi{\arabic{enumi}.}
\tightlist
\item
  \textbf{Trial \(\alpha_U\)} \(\rightarrow\) run couplings to scale \(\mu\).
\item
  \textbf{Compute} \(v = E_{\text{cell}} \cdot \exp(-2\pi/(\beta_{\text{EW}} \cdot \alpha_U))\) with \(\beta_{\text{EW}} = N_c + 1 = 4\).
\item
  \textbf{Compute} tree-level \(m_Z(\mu) = \frac{1}{2} v \sqrt{g_2^2 + g_Y^2}\) where \(g_Y = \sqrt{4\pi \cdot \frac{3}{5}\alpha_1}\).
\item
  \textbf{Find} \(\mu^* = m_Z(\mu^*)\) (self-consistent fixed point).
\item
  \textbf{Evaluate} the pixel residual: \(\bar{\ell}_{\text{SU(2)}}(t_2) + \bar{\ell}_{\text{SU(3)}}(t_3) - P/4\).
\item
  \textbf{Bisect} on \(\alpha_U\) until the pixel residual vanishes.
\end{enumerate}

This yields:

\[\alpha_U \approx 0.04112, \quad \alpha_U^{-1} \approx 24.32\]

\begin{center}\rule{0.5\linewidth}{0.5pt}\end{center}

\subsection{4. Stage 3: Electroweak Observables}\label{4-stage-3-electroweak-observables}

With \(\alpha_U\) and \(M_U\) determined, the gauge couplings at \(\mu = m_{Z,\text{run}}\) are fixed:

{\def\LTcaptype{none} % do not increment counter
\begin{longtable}[]{@{}>{\RaggedRight\arraybackslash}p{0.320\textwidth}>{\RaggedRight\arraybackslash}p{0.320\textwidth}>{\RaggedRight\arraybackslash}p{0.320\textwidth}@{}}
\toprule\noalign{}
Quantity & Formula & Predicted Value \\
\midrule\noalign{}
\endhead
\bottomrule\noalign{}
\endlastfoot
\(m_{Z,\text{run}}\) & Self-consistent fixed point & 91.652 GeV \\
\(v\) (Higgs VEV) & \(E_{\text{cell}} \cdot e^{-2\pi/(\beta_{\text{EW}} \alpha_U)}\) & 246.77 GeV \\
\(\alpha_1(m_{Z,\text{run}})\) & \([\alpha_U^{-1} + \frac{b_1}{2\pi}\ln(M_U/m_{Z,\text{run}})]^{-1}\) & 0.01696 \\
\(\alpha_2(m_{Z,\text{run}})\) & \([\alpha_U^{-1} + \frac{b_2}{2\pi}\ln(M_U/m_{Z,\text{run}})]^{-1}\) & 0.03384 \\
\(\alpha_3(m_{Z,\text{run}})\) & \([\alpha_U^{-1} + \frac{b_3}{2\pi}\ln(M_U/m_{Z,\text{run}})]^{-1}\) & 0.1183 \\
\end{longtable}
}

From these:

\[\alpha_{\text{em}} = \left(\frac{1}{\alpha_2} + \frac{1}{\frac{3}{5}\alpha_1}\right)^{-1}, \quad \sin^2\theta_W = \frac{\alpha_{\text{em}}}{\alpha_2}\]

We use the electroweak normalization
\[
\alpha_1=\frac53\alpha_Y,
\qquad
g_Y^2=4\pi\alpha_Y=4\pi\frac35\alpha_1.
\]

{\def\LTcaptype{none} % do not increment counter
\begin{longtable}[]{@{}>{\RaggedRight\arraybackslash}p{0.320\textwidth}>{\RaggedRight\arraybackslash}p{0.320\textwidth}>{\RaggedRight\arraybackslash}p{0.320\textwidth}@{}}
\toprule\noalign{}
Observable & Predicted & PDG Value \\
\midrule\noalign{}
\endhead
\bottomrule\noalign{}
\endlastfoot
\(\alpha_{\text{em}}^{-1}(m_{Z,\text{run}})\) & 128.31 & 127.952 \\
\(\sin^2\theta_W(m_{Z,\text{run}})\) & 0.2307 & 0.23122 \\
\(\alpha_s(m_{Z,\text{run}})\) & 0.1183 & 0.1179 \
\end{longtable}
}

\subsubsection{4.1 Z Boson Pole Mass}\label{41-z-boson-pole-mass}

The tree-level \(m_{Z,\text{run}} = 91.65\) GeV is the running mass at its own scale. The physical pole mass includes the custodial correction from top-quark loops:

\[\Delta\rho = \frac{3y_t^2}{32\pi^2}\]

For the stage-3 numerical estimate we use \(y_t \approx 1\), giving

\[\Delta\rho_{\text{stage-3}} \approx \frac{3}{32\pi^2} \approx 0.009499\]

\[M_{Z,\text{pole}} = \frac{m_{Z,\text{run}}}{\sqrt{1 + \Delta\rho}} \approx 91.220 \text{ GeV}\]

\textbf{Why the stage-3 estimate uses \(\Delta\rho \approx 3/(32\pi^2)\).} The general one-loop custodial-symmetry-breaking contribution is \(\Delta\rho = 3y_t^2/(32\pi^2)\). The displayed stage-3 number sets \(y_t \approx 1\), requiring no knowledge of the actual top mass , only that the top Yukawa is order-one (the least-suppressed channel in the \(\mathbb{Z}_6\) texture).

\subsubsection{4.2 W Boson Mass}\label{42-w-boson-mass}

At tree level:

\[m_W = \frac{1}{2} v \cdot g_2 = \frac{1}{2} v \sqrt{4\pi\alpha_2} \approx 80.39 \text{ GeV}\]

\begin{center}\rule{0.5\linewidth}{0.5pt}\end{center}

\subsection{5. Stage 4: Critical Surface , Higgs and Top}\label{5-stage-4-critical-surface--higgs-and-top}

\subsubsection{5.1 The Critical Surface Constraint}\label{51-the-critical-surface-constraint}

Refinement stability (§6.3, §6.22 of Part I of this manuscript) pushes the Higgs quartic coupling to a marginal stability point at the unification scale:

\[\lambda(M_U) = 0, \quad \beta_\lambda(M_U) = 0\]

Setting \(\beta_\lambda = 0\) at one loop with \(\lambda = 0\) fixes the top Yukawa boundary condition:

\[y_t(M_U) = \left[\frac{1}{16}\left(2g_2^4 + (g_2^2 + g_Y^2)^2\right)\right]^{1/4}\]

where \(g_Y(M_U)\) and \(g_2(M_U)\) are obtained inside the synchronized D11 transport by running the OPH-derived couplings from \(m_Z\) to a UV synchronization scale \(\mu_{\mathrm{sync}}\approx 6.8\times 10^{11}\,\mathrm{GeV}\) with SM coefficients and from \(\mu_{\mathrm{sync}}\) to \(M_U\) with the same MSSM-like coefficients used in D10. The raw all-SM appendix flow by itself lands far lower (\(m_t^{\mathrm{pole}}\approx 164.2\,\mathrm{GeV}\), \(m_H\approx 115.2\,\mathrm{GeV}\)); the synchronized D11 numbers quoted below are the UV-synchronized transport plus only order-one-GeV low-scale matching corrections.

\subsubsection{5.2 RG Evolution to Low Scales}\label{52-rg-evolution-to-low-scales}

The coupled system \((y_t, \lambda, g_Y, g_2, g_3)\) is integrated from \(M_U\) down to \(\mu \sim m_t\) using RK4. The synchronized D11 implementation uses a piecewise one-loop gauge transport:

\[\alpha_i^{-1}(\mu) = \alpha_i^{-1}(m_Z) - \frac{b_i^{\text{SM}}}{2\pi}\ln(\mu/m_Z), \qquad m_Z \le \mu \le \mu_{\mathrm{sync}}\]

and
\[\alpha_i^{-1}(\mu) = \alpha_i^{-1}(\mu_{\mathrm{sync}}) - \frac{b_i^{\text{UV}}}{2\pi}\ln(\mu/\mu_{\mathrm{sync}}), \qquad \mu_{\mathrm{sync}} \le \mu \le M_U,\]
with \(b^{\text{UV}}=(33/5,1,-3)\) and \(\mu_{\mathrm{sync}}\approx 6.8\times 10^{11}\,\mathrm{GeV}\). This Stage-4 section is the authoritative synchronized D11 implementation used for the quoted Higgs/top numbers. The compact appendix records only the low-scale SM part of the one-loop system; the numerical outputs below also use the UV synchronization, low-scale extraction, and pole-mass conventions fixed here in the synchronized supplement-backed branch.

\subsubsection{5.3 Top Mass}\label{53-top-mass}

The UV-synchronized running top mass is determined self-consistently: iterate \(\mu_{\text{guess}} = y_t(\mu) \cdot v/\sqrt{2}\) until convergence:

\[m_t^{\overline{\text{MS}}}(m_t) \approx 160.6 \text{ GeV}\]

The same UV-synchronized core gives
\[m_t^{\text{pole,core}} \approx 170.3 \text{ GeV}\]
before the final low-scale matching adjustment. The pole-mass part of that core uses the standard 3-loop QCD relation:

\[\frac{m_t^{\text{pole}}}{m_t^{\overline{\text{MS}}}} = 1 + \frac{4}{3}\frac{\alpha_s}{\pi} + K_2 \left(\frac{\alpha_s}{\pi}\right)^2 + K_3 \left(\frac{\alpha_s}{\pi}\right)^3\]

with \(K_2 = 13.44 - 1.04 n_l\) and \(K_3 = 190.6 - 26.7 n_l + 0.65 n_l^2\) (\(n_l = 5\)). The displayed synchronized D11 output then adds a residual \(+0.84\) GeV matching correction:

\[m_t^{\text{pole}} \approx 171.1 \text{ GeV}\]

\subsubsection{5.4 Higgs Mass}\label{54-higgs-mass}

From \(\lambda(m_t)\) obtained by the UV-synchronized RG integration one finds a core value
\[m_H^{\text{core}} \approx 126.6 \text{ GeV}\]
and after the residual \(-0.12\) GeV Higgs matching shift used in the synchronized implementation:

\[m_H = \sqrt{2\lambda(m_t)} \cdot v \approx 126.5 \text{ GeV}\]

\begin{center}\rule{0.5\linewidth}{0.5pt}\end{center}

\subsection{6. Stage 5: Discrete Spectrum , Quarks and Leptons}\label{6-stage-5-discrete-spectrum--quarks-and-leptons}

\subsubsection{\texorpdfstring{6.1 The \(\mathbb{Z}_6\) Defect Parameter}{6.1 The \textbackslash mathbb\{Z\}\_6 Defect Parameter}}\label{61-the-mathbbz_6-defect-parameter}

The SM global gauge group is \((SU(3) \times SU(2) \times U(1))/\mathbb{Z}_6\). The entropy deficit from the \(\mathbb{Z}_6\) quotient yields the universal suppression factor:

\[\varepsilon = e^{-\ln 6} = \frac{1}{6}\]

This is the base of the Froggatt-Nielsen texture: each unit of \(\mathbb{Z}_6\) defect insertion suppresses a Yukawa coupling by a factor of 6.

\subsubsection{6.2 Integer Exponents (Algorithmically Derived)}\label{62-integer-exponents-algorithmically-derived}

The diagonal mass exponents are determined by the minimal SU(3) hierarchy structure with \(N_c = 3\) colors and \(N_g = 3\) generations:

\textbf{Up-type quarks}: \(n_u = (2N_c, N_c, 0) = (6, 3, 0)\)

\textbf{Down-type quarks}: \(n_d = (2N_c, N_c+1, N_c-1) = (6, 4, 2)\)

\textbf{Charged leptons}: derived from the Koide phase.

\subsubsection{6.3 The Koide Phase}\label{63-the-koide-phase}

The charged lepton masses use a circulant ansatz on generation space with Koide ratio \(Q = 2/3\) (from \(\mathbb{Z}_3\) mode balance). The holonomy phase is:

\[\delta = \frac{\beta_{\text{EW}}}{2 N_c N_g} = \frac{N_c + 1}{2 N_c N_g} = \frac{4}{18} = \frac{2}{9}\]

The Koide roots are:

\[r_k = 1 + \sqrt{2}\cos\left(\frac{2}{9} + \frac{2\pi k}{3}\right), \quad k = 0, 1, 2\]

sorted in ascending order. The lepton exponents \((n_e, n_\mu, n_\tau)\) are selected by matching ratios \(r_i^2/r_j^2 \approx \varepsilon^{n_i - n_j}\), yielding \((n_e, n_\mu, n_\tau) = (7, 4, 3)\).

\subsubsection{6.4 Diagonal Quark Masses}\label{64-diagonal-quark-masses}

\[m_{u_i} = \frac{v}{\sqrt{2}} \varepsilon^{n_{u,i}}, \quad m_{d_i} = \frac{v}{\sqrt{2}} \varepsilon^{n_{d,i}}\]

\subsubsection{6.5 CKM Mixing}\label{65-ckm-mixing}

The CKM matrix is parameterized by \(\varepsilon\):

\[s_{12} = \varepsilon = 1/6, \quad s_{23} = \varepsilon^2 = 1/36, \quad s_{13} = \varepsilon^3 = 1/216\]

Physical down-type masses are the singular values of \(V_{\text{CKM}} \cdot \text{diag}(m_d, m_s, m_b)\).

\subsubsection{6.6 Charged Lepton Masses}\label{66-charged-lepton-masses}

A scale factor \(S\) is determined from the exponents \(n_e = (7, 4, 3)\) and the Koide roots by demanding internal consistency:

\[\ln S_0 = -\frac{1}{3}\sum_{i=1}^3 \ln\left(\frac{r_i^2 \sqrt{2}}{v} \cdot 6^{n_{e,i}}\right)\]

Then \(S = S_0 \cdot 2^{1/6}\) and:

\[m_e = S \cdot r_1^2, \quad m_\mu = S \cdot r_2^2, \quad m_\tau = S \cdot r_3^2\]

\begin{center}\rule{0.5\linewidth}{0.5pt}\end{center}

\subsection{7. Stage 6: QCD Scale and Hadron Masses}\label{7-stage-6-qcd-scale-and-hadron-masses}

\subsubsection{\texorpdfstring{7.1 \(\Lambda_{\overline{\text{MS}}}\) from \(\alpha_s(m_Z)\)}{7.1 \textbackslash Lambda\_\{\textbackslash overline\{\textbackslash text\{MS\}\}\} from \textbackslash alpha\_s(m\_Z)}}\label{71-lambda_overlinetextms-from-alpha_sm_z}

Given the OPH-derived \(\alpha_s(m_Z) \approx 0.1183\) (a calibration-sector output) and the quark thresholds (\(m_b\), \(m_c\) from Stage 5), the QCD scale is extracted using the standard 4-loop MSbar definition:

\[\ln\frac{\mu^2}{\Lambda^2} = \frac{1}{\beta_0 a} + \frac{\beta_1}{\beta_0^2}\ln(\beta_0 a) + \int_0^a dx\left[\frac{1}{\beta(x)} + \frac{1}{\beta_0 x^2} - \frac{\beta_1}{\beta_0^2 x}\right]\]

where \(a = \alpha_s/(4\pi)\) and \(\beta_0, \beta_1, \beta_2, \beta_3\) are the 4-loop MSbar coefficients.

Stepping down across flavor thresholds:

\[\alpha_s^{(n_f-1)}(\mu_{\text{th}}) = \alpha_s^{(n_f)}(\mu_{\text{th}}) \quad \text{at } \mu_{\text{th}} = m_b, m_c\]

yields:

{\def\LTcaptype{none} % do not increment counter
\begin{longtable}[]{@{}>{\RaggedRight\arraybackslash}p{0.480\textwidth}>{\RaggedRight\arraybackslash}p{0.480\textwidth}@{}}
\toprule\noalign{}
\(n_f\) & \(\Lambda_{\overline{\text{MS}}}^{(n_f)}\) (GeV) \\
\midrule\noalign{}
\endhead
\bottomrule\noalign{}
\endlastfoot
5 & 0.211 \\
4 & 0.288 \\
3 & 0.322 \\
\end{longtable}
}

\subsubsection{7.2 Hadron Masses}\label{72-hadron-masses}

Hadron masses require a nonperturbative computation of dimensionless ratios \(C_X = m_X / \Lambda^{(3)}\). These are computed by an internal lattice QCD collar calculation:

\begin{enumerate}
\def\labelenumi{\arabic{enumi}.}
\tightlist
\item
  \textbf{Wilson SU(3) gauge} (quenched) with Metropolis updates.
\item
  \textbf{Wilson valence quarks} at multiple \(\kappa\) values for chiral extrapolation.
\item
  \textbf{Gradient-flow coupling} for input-free \(a\Lambda\) determination.
\item
  \textbf{Two-point Richardson extrapolation} for the continuum limit.
\end{enumerate}

Then: \(m_X = C_X \cdot \Lambda_{\overline{\text{MS}}}^{(3)}\).

The lattice computation uses small volumes (\(L = 2\)--\(6\)) in the quenched approximation. Systematic uncertainties from quenching (\textasciitilde10\%), finite volume, and limited statistics dominate the error budget for hadron masses.

\begin{center}\rule{0.5\linewidth}{0.5pt}\end{center}

\subsection{8. Stage 7: Neutrino Masses}\label{8-stage-7-neutrino-masses}

\subsubsection{8.1 The Capacity Model}\label{81-the-capacity-model}

The second OPH input , the screen capacity \(\log(\dim\mathcal{H}) \sim 10^{122}\) , enters through the cosmological constant:

\[\rho_\Lambda = \frac{3}{8} \frac{M_P^4}{\log(\dim\mathcal{H})}\]

The neutrino mass scale is anchored at \(\rho_\Lambda^{1/4}\):

\[m_{\nu_3} = \rho_\Lambda^{1/4} \approx 3.0 \times 10^{-12} \text{ GeV} \approx 3.0 \text{ meV}\]

The hierarchy uses the same \(\varepsilon = 1/6\) suppression:

\[m_{\nu_2} = \varepsilon \cdot m_{\nu_3} \approx 0.50 \text{ meV}, \quad m_{\nu_1} = \varepsilon^2 \cdot m_{\nu_3} \approx 0.084 \text{ meV}\]

This yields \(\Delta m_{32}^2 \approx 9.1 \times 10^{-24} \text{ GeV}^2\) and \(\Delta m_{21}^2 \approx 2.5 \times 10^{-25} \text{ GeV}^2\). These values remain well below current direct-detection bounds, but they do not yet reproduce the observed atmospheric and solar oscillation splittings.

\begin{center}\rule{0.5\linewidth}{0.5pt}\end{center}

\subsection{9. Massless Particles}\label{9-massless-particles}

The following particles are predicted to be \textbf{exactly massless} by symmetry:

{\def\LTcaptype{none} % do not increment counter
\begin{longtable}[]{@{}>{\RaggedRight\arraybackslash}p{0.480\textwidth}>{\RaggedRight\arraybackslash}p{0.480\textwidth}@{}}
\toprule\noalign{}
Particle & Reason \\
\midrule\noalign{}
\endhead
\bottomrule\noalign{}
\endlastfoot
Photon & Unbroken \(U(1)_{\text{em}}\) gauge invariance forbids a mass term \\
Gluons (8) & Unbroken \(SU(3)_c\) gauge invariance forbids a mass term \\
Graviton & Diffeomorphism invariance (emergent from entanglement equilibrium) forbids a mass term \\
\end{longtable}
}

These are \textbf{symmetry-protected zeros}, not accidental cancellations. The gauge invariances are derived from the gluing structure (§6.1 of Part I of this manuscript) and entanglement equilibrium (§5 of Part I of this manuscript).

\begin{center}\rule{0.5\linewidth}{0.5pt}\end{center}

\subsection{10. Complete Spectrum: OPH Outputs vs PDG}\label{10-complete-spectrum-oph-outputs-vs-pdg}

All outputs below use inputs \(P = 1.63094\) and \(\log(\dim\mathcal{H}) = 10^{122}\) only. PDG values are from the 2024/2025 edition. Each output is classified as \textbf{consistency check} (Tier~1, from calibration sector), \textbf{independent quantitative} (Tier~2), \textbf{phenomenological continuation / downstream branch} (Tier~3), or \textbf{structural} (Tier~4, parameter-free). See §1B for the full classification.

\subsubsection{10.1 Gauge Bosons and Higgs}\label{101-gauge-bosons-and-higgs}

{\def\LTcaptype{none} % do not increment counter
\begin{longtable}[]{@{}>{\RaggedRight\arraybackslash}p{0.160\textwidth}>{\RaggedLeft\arraybackslash}p{0.160\textwidth}>{\RaggedLeft\arraybackslash}p{0.160\textwidth}>{\RaggedLeft\arraybackslash}p{0.160\textwidth}>{\RaggedRight\arraybackslash}p{0.160\textwidth}>{\RaggedRight\arraybackslash}p{0.160\textwidth}@{}}
\toprule\noalign{}
Particle & OPH (GeV) & PDG (GeV) & Rel. Error & Tier & Error Source \\
\midrule\noalign{}
\endhead
\bottomrule\noalign{}
\endlastfoot
\(\gamma\) & 0 & 0 & exact & Structural & Symmetry-protected; no error expected \\
\(g\) (gluon) & 0 & 0 & exact & Structural & Symmetry-protected; no error expected \\
\(W\) boson & 80.386 & 80.377 ± 0.012 & +0.012\% & Consistency & Flows from calibrated \(\alpha_U\) and \(v\). Tree-level only \\
\(Z\) pole & 91.220 & 91.188 ± 0.002 & +0.035\% & Consistency & Flows from calibrated couplings. Simplified \(\Delta\rho\) \\
\(H\) (Higgs) & 126.48 & 125.20 ± 0.11 & +1.02\% & \textbf{Prediction} & D11 supplement reconstruction: UV-synchronized SM/MSSM transport plus residual low-scale matching after \(\lambda = \beta_\lambda = 0\) \\
\end{longtable}
}

\subsubsection{10.2 Charged Leptons}\label{102-charged-leptons}

{\def\LTcaptype{none} % do not increment counter
\begin{longtable}[]{@{}>{\RaggedRight\arraybackslash}p{0.160\textwidth}>{\RaggedLeft\arraybackslash}p{0.160\textwidth}>{\RaggedLeft\arraybackslash}p{0.160\textwidth}>{\RaggedLeft\arraybackslash}p{0.160\textwidth}>{\RaggedRight\arraybackslash}p{0.160\textwidth}>{\RaggedRight\arraybackslash}p{0.160\textwidth}@{}}
\toprule\noalign{}
Particle & OPH (GeV) & PDG (GeV) & Rel. Error & Tier & Error Source \\
\midrule\noalign{}
\endhead
\bottomrule\noalign{}
\endlastfoot
Electron & 5.109 × 10⁻⁴ & 5.110 × 10⁻⁴ & −0.023\% & \textbf{Prediction} & Koide structure (\(Q=2/3\), \(\delta=2/9\)) independent of calibration; residual from QED running \\
Muon & 0.10564 & 0.10566 & −0.022\% & \textbf{Prediction} & Same Koide structure; QED running corrections \textasciitilde0.02\%, matching the offset \\
Tau & 1.7766 & 1.7769 ± 0.09 & −0.020\% & \textbf{Prediction} & Same origin. Uniform sign shows common scale-factor effect from QED running \\
\end{longtable}
}

\subsubsection{10.3 Neutrinos}\label{103-neutrinos}

{\def\LTcaptype{none} % do not increment counter
\begin{longtable}[]{@{}>{\RaggedRight\arraybackslash}p{0.192\textwidth}>{\RaggedLeft\arraybackslash}p{0.192\textwidth}>{\RaggedLeft\arraybackslash}p{0.192\textwidth}>{\RaggedRight\arraybackslash}p{0.192\textwidth}>{\RaggedRight\arraybackslash}p{0.192\textwidth}@{}}
\toprule\noalign{}
Particle & OPH (GeV) & Experiment & Tier & Error Source \\
\midrule\noalign{}
\endhead
\bottomrule\noalign{}
\endlastfoot
\(\nu_e\) & 8.39 × 10⁻¹⁴ & \(< 8 \times 10^{-10}\) (KATRIN) & \textbf{Prediction} & Uses second input (\(\log\dim\mathcal{H}\)), independent of \(P\). Consistent with all bounds \\
\(\nu_\mu\) & 5.04 × 10⁻¹³ & \(< 8 \times 10^{-10}\) (KATRIN) & \textbf{Prediction} & Same capacity branch; below direct bounds, but not yet matched to the solar oscillation scale \\
\(\nu_\tau\) & 3.02 × 10⁻¹² & \(< 8 \times 10^{-10}\) (KATRIN) & \textbf{Prediction} & Anchored at \(\rho_\Lambda^{1/4}\); below direct bounds, but not yet matched to the atmospheric oscillation scale \\
\end{longtable}
}

All predictions are well below the current experimental upper bound of \(\sim 0.8\) eV \(\approx 8 \times 10^{-10}\) GeV. At present this branch is better described as a low-scale capacity estimate than as a quantitative fit to oscillation data.

\subsubsection{10.4 Quarks}\label{104-quarks}

{\def\LTcaptype{none} % do not increment counter
\begin{longtable}[]{@{}>{\RaggedRight\arraybackslash}p{0.160\textwidth}>{\RaggedLeft\arraybackslash}p{0.160\textwidth}>{\RaggedLeft\arraybackslash}p{0.160\textwidth}>{\RaggedLeft\arraybackslash}p{0.160\textwidth}>{\RaggedRight\arraybackslash}p{0.160\textwidth}>{\RaggedRight\arraybackslash}p{0.160\textwidth}@{}}
\toprule\noalign{}
Particle & OPH (GeV) & PDG (GeV) & Rel. Error & Tier & Error Source \\
\midrule\noalign{}
\endhead
\bottomrule\noalign{}
\endlastfoot
up & 3.74 × 10⁻³ & 2.16 × 10⁻³ & +73\% & \textbf{Prediction} & \textbf{\(u\)-\(d\) degeneracy}: both get exponent \(n=6\). Subleading effects not included \\
down & 3.74 × 10⁻³ & 4.70 × 10⁻³ & −20\% & \textbf{Prediction} & Same degeneracy. Geometric mean \(\sqrt{m_u m_d} \approx 3.2\) MeV closer to 3.7 MeV \\
strange & 0.1346 & 0.0935 & +44\% & \textbf{Prediction} & No scheme matching; missing Clebsch factors and CKM rotation \\
charm & 0.808 & 1.273 & −37\% & \textbf{Prediction} & Missing RG running from \(v\) to \(m_c\) (\textasciitilde factor 1.3) and Clebsch coefficient \\
bottom & 4.847 & 4.183 & +16\% & \textbf{Prediction} & Running from \(v\) to \(m_b\) reduces mass by \textasciitilde15\%, closing most of the gap \\
top (texture) & 174.5 & 172.6 & +1.1\% & \textbf{Prediction} & Exponent \(n_t = 0\) (unsuppressed Yukawa); small error from \(v\) offset \\
top (crit. surf.) & 171.1 & 172.6 & −0.87\% & \textbf{Prediction} & D11 supplement reconstruction: UV-synchronized SM/MSSM transport plus residual pole-mass matching after \(\lambda = \beta_\lambda = 0\) \\
\end{longtable}
}

\subsubsection{\texorpdfstring{10.5 Gauge Couplings at \(m_{Z,\text{run}}\)}{10.5 Gauge Couplings at m\_Zrun}}\label{105-gauge-couplings-at-m_z}

These are \textbf{calibration-sector consistency checks} (Tier~1). The gauge couplings at \(m_{Z,\text{run}}\) are the fixed-point values used in the entropy-matching condition and self-consistent closure procedure. Their agreement with PDG values demonstrates internal consistency, not independent predictive power.

{\def\LTcaptype{none} % do not increment counter
\begin{longtable}[]{@{}>{\RaggedRight\arraybackslash}p{0.18\textwidth}>{\RaggedLeft\arraybackslash}p{0.14\textwidth}>{\RaggedLeft\arraybackslash}p{0.18\textwidth}>{\RaggedLeft\arraybackslash}p{0.12\textwidth}>{\RaggedRight\arraybackslash}p{0.10\textwidth}>{\RaggedRight\arraybackslash}p{0.18\textwidth}@{}}
\toprule\noalign{}
Quantity & OPH & PDG & Rel. Error & Tier & Error Source \\
\midrule\noalign{}
\endhead
\bottomrule\noalign{}
\endlastfoot
\(\alpha_{\text{em}}^{-1}(m_Z)\) & 128.31 & 127.952 ± 0.009 & +0.28\% & Consistency & Flows from calibrated \(\alpha_U\) \\
\(\sin^2\theta_W(m_Z)\) & 0.2307 & 0.23122 ± 0.00004 & −0.21\% & Consistency & Most sensitive to threshold scale \\
\(\alpha_s(m_Z)\) & 0.1183 & 0.1179 ± 0.0009 & +0.37\% & Consistency & Within 0.5\(\sigma\) of PDG \\
\end{longtable}
}

\subsubsection{10.6 Derived Scales}\label{106-derived-scales}

{\def\LTcaptype{none} % do not increment counter
\begin{longtable}[]{@{}>{\RaggedRight\arraybackslash}p{0.14\textwidth}>{\RaggedLeft\arraybackslash}p{0.13\textwidth}>{\RaggedLeft\arraybackslash}p{0.13\textwidth}>{\RaggedLeft\arraybackslash}p{0.10\textwidth}>{\RaggedRight\arraybackslash}p{0.13\textwidth}>{\RaggedRight\arraybackslash}p{0.27\textwidth}@{}}
\toprule\noalign{}
Quantity & OPH Value & Reference & Agreement & Tier & Error Source \\
\midrule\noalign{}
\endhead
\bottomrule\noalign{}
\endlastfoot
\(v\) (Higgs VEV) & 246.77 GeV & 246.22 GeV & +0.22\% & Consistency & Flows from calibrated \(\alpha_U\) via transmutation formula \\
\(M_U\) & 2.47 × 10¹⁶ GeV & --- & --- & Consistency & No direct measurement; consistent with proton stability bounds \\
\(\alpha_U^{-1}\) & 24.32 & --- & --- & Consistency & Intermediate variable in calibration procedure \\
\(\Lambda_{\overline{\text{MS}}}^{(3)}\) & 0.322 GeV & \textasciitilde0.332 GeV & −3.0\% & Consistency & Propagated from \(\alpha_s\) (calibration sector); 0.4\% in \(\alpha_s\) \(\to\) \textasciitilde3\% in \(\Lambda\) \\
\end{longtable}
}

\subsubsection{10.7 Hadron Masses}\label{107-hadron-masses}

Hadron masses follow from \(m_X = C_X \cdot \Lambda_{\overline{\text{MS}}}^{(3)}\), where \(C_X\) is a dimensionless ratio computed via lattice QCD. The derivation chain is complete: \(P \to \alpha_s(m_Z) \to \Lambda^{(3)} \to m_X\). Current lattice estimates use small volumes and the quenched approximation, so systematic uncertainties are large (\textasciitilde10--20\%).

{\def\LTcaptype{none} % do not increment counter
\begin{longtable}[]{@{}>{\RaggedRight\arraybackslash}p{0.240\textwidth}>{\RaggedLeft\arraybackslash}p{0.240\textwidth}>{\RaggedRight\arraybackslash}p{0.240\textwidth}>{\RaggedRight\arraybackslash}p{0.240\textwidth}@{}}
\toprule\noalign{}
Particle & PDG (GeV) & Formula & Dominant uncertainty \\
\midrule\noalign{}
\endhead
\bottomrule\noalign{}
\endlastfoot
Proton & 0.93827 & \(C_p \cdot \Lambda^{(3)}\), \(C_p \approx 2.9\) & Quenching (\textasciitilde10\%), finite volume \\
Neutron & 0.93957 & \(\approx m_p\) in the isospin limit & Isospin splitting requires unquenched QCD+QED \\
\(\pi^\pm\) & 0.13957 & \(C_\pi \cdot \Lambda^{(3)}\) & Chiral extrapolation (\(m_\pi^2 \propto m_q\)) \\
\(\pi^0\) & 0.13498 & \(\approx m_{\pi^\pm}\) & EM splitting \\
\(K^\pm\) & 0.4937 & \(C_K \cdot \Lambda^{(3)}\) & Strange-quark valence mass \\
\(K^0\) & 0.4976 & \(\approx m_{K^\pm}\) & EM splitting \\
\(\Lambda\) baryon & 1.1157 & \(C_\Lambda \cdot \Lambda^{(3)}\) & Strange-quark baryon correlator \\
\end{longtable}
}

\begin{center}\rule{0.5\linewidth}{0.5pt}\end{center}

\subsection{11. Analysis of Discrepancies}\label{11-analysis-of-discrepancies}

\subsubsection{11.1 Sub-Permille Outputs (\textless{} 0.1\%)}\label{111-sub-permille-outputs--01}

\textbf{W boson} (+0.012\%), \textbf{Z boson} (+0.035\%), \textbf{electron} (−0.023\%), \textbf{muon} (−0.022\%), \textbf{tau} (−0.020\%):

These five outputs are at the sub-permille level, but their epistemic status differs. The gauge boson masses (\(W\), \(Z\)) are \textbf{consistency checks}: they flow directly from the calibrated gauge couplings and \(v\), so their agreement is expected by construction. The charged leptons (\(e\), \(\mu\), \(\tau\)) are \textbf{Tier~3 phenomenological continuations}: the Koide structure (\(Q = 2/3\), \(\delta = 2/9\)) is algebraically independent of the gauge-coupling calibration, but the final exponent assignment uses an additional discrete-selection step.

\subsubsection{11.2 Percent-Level Outputs (0.1\% -- 2\%)}\label{112-percent-level-outputs-01--2}

\textbf{Higgs} (+1.0\%), \textbf{top pole} (−0.87\%), \textbf{\(\alpha_s\)} (+0.37\%), \textbf{\(\sin^2\theta_W\)} (−0.21\%), \textbf{\(\alpha_{\text{em}}^{-1}\)} (+0.28\%):

The epistemic status differs within this group. The \textbf{gauge couplings} (\(\alpha_s\), \(\sin^2\theta_W\), \(\alpha_{\text{em}}^{-1}\)) are \textbf{consistency checks} --- they are the data used to calibrate \(P\), so their agreement demonstrates that the self-consistent fixed-point procedure works, not that the framework predicts them.

The \textbf{Higgs and top masses} are \textbf{Tier~2 independent quantitative outputs}: they come from the critical surface constraint (\(\lambda(M_U) = 0\), \(\beta_\lambda(M_U) = 0\)) together with the synchronized D11 transport (SM running below \(\mu_{\mathrm{sync}}\), MSSM-like running above, and only order-one-GeV residual matching). This remains an independent boundary condition not used in calibrating \(P\). The \textasciitilde1\% error is consistent with the expected size of 2-loop corrections and threshold effects.

\subsubsection{11.3 Qualitative Agreement Only: Quark Masses (16\% -- 73\%)}\label{113-qualitative-agreement-only-quark-masses-16--73}

The quark masses from the \(\varepsilon = 1/6\) texture show large deviations. The reasons are well-understood:

\begin{enumerate}
\def\labelenumi{\arabic{enumi}.}
\item
  \textbf{\(u\)-\(d\) degeneracy}: The texture assigns identical exponents \(n_u = n_d = 6\) to the up and down quarks, predicting \(m_u = m_d\). In reality, \(m_u/m_d \approx 0.46\). Breaking this degeneracy requires isospin-violating effects beyond the minimal \(\mathbb{Z}_6\) texture (e.g., subleading defect insertions or electromagnetic corrections).
\item
  \textbf{Scheme and scale mismatch}: The PDG quark masses are \(\overline{\text{MS}}\) running masses at scale \(\mu = 2\) GeV (for light quarks) or \(\mu = m_q\) (for heavy quarks). The OPH texture produces "Yukawa-scale" masses at \(\mu \sim v\) with no scheme matching applied. For the charm quark, this can easily account for a factor of \textasciitilde1.5.
\item
  \textbf{Order-one Clebsch-Gordan coefficients}: The texture \(y_f = c_f \cdot \varepsilon^{n_f}\) has residual coefficients \(c_f\) that are undetermined (expected to be \(\mathcal{O}(1)\), actually ranging from 0.6 to 2.2). These encode CKM rotation effects, RG running from \(M_U\) to \(m_Z\), and overlap matrix elements that the minimal integer-exponent ansatz does not resolve.
\item
  \textbf{Missing threshold corrections}: No QCD or electroweak threshold corrections are applied at flavor thresholds. These are typically 5--20\% effects for the lighter quarks.
\end{enumerate}

\textbf{The correct interpretation}: The texture correctly captures the \emph{hierarchy} (the fact that \(m_t/m_u \sim 10^5\) arises from integer exponents in base 6), and it does not aim for precision at the individual-mass level. The base-6 logarithms of all nine charged-fermion Yukawas land within \textasciitilde0.3 of integers , this is the core prediction. The order-one residuals are where scheme matching and subleading effects live.

\subsubsection{11.4 Hadron Masses}\label{114-hadron-masses}

The hadron mass predictions follow from the dimensionless ratios \(C_X = m_X/\Lambda^{(3)}\), computed via lattice QCD. The internal lattice code (\texttt{oph\_lattice\_su3\_quenched\_v5.py}) implements a quenched Wilson-gauge SU(3) lattice with Wilson valence quarks, gradient-flow scale setting, and Richardson continuum extrapolation , all without PDG inputs. The dominant systematic uncertainties are:

\begin{itemize}
\tightlist
\item
  \textbf{Quenching error} (\textasciitilde10\%): The gauge field is quenched (\(n_f = 0\)), while physical QCD has \(n_f = 2+1\) light dynamical flavors.
\item
  \textbf{Volume effects}: Lattice volumes (\(L = 2\)--\(6\)) give finite-volume corrections.
\item
  \textbf{Statistics}: Current runs use \(\mathcal{O}(10)\) configurations; scaling to \(\mathcal{O}(10^3)\) would reduce statistical errors.
\end{itemize}

These are standard lattice QCD systematics, not limitations of the OPH framework. The derivation chain \(P \to \alpha_s \to \Lambda^{(3)} \to m_X\) is complete; improving the lattice computation improves the numerical precision of the hadron predictions.

\subsubsection{11.5 Neutrinos: No Direct Comparison Available}\label{115-neutrinos-no-direct-comparison-available}

The predicted neutrino masses (\(\sim 0.08\)--\(3\) meV) are below current direct-detection sensitivity (KATRIN: \(m_\beta < 0.8\) eV) and well below current cosmological mass-sum bounds. The predicted \(\sum m_\nu \approx 3.6\) meV is well below the cosmological bound \(\sum m_\nu < 0.12\) eV (Planck 2018), but the simple capacity-branch hierarchy does not yet reproduce the observed oscillation splittings quantitatively.

\begin{center}\rule{0.5\linewidth}{0.5pt}\end{center}

\subsection{12. Reproduction Instructions}\label{12-reproduction-instructions}

All computations can be reproduced from the code in \href{../code/particles/}{\texttt{code/particles/}}.

\subsubsection{12.1 Prerequisites}\label{121-prerequisites}

\begin{verbatim}
pip install numpy
\end{verbatim}

No other dependencies are needed. The \texttt{pdg} and \texttt{pandas} packages are only required for \path{tools/fetch_pdg_data.py} (the PDG data fetcher, used for comparison only).

\subsubsection{12.2 Code Files}\label{122-code-files}

{\def\LTcaptype{none} % do not increment counter
\begin{longtable}[]{@{}>{\RaggedRight\arraybackslash}p{0.390\textwidth}>{\RaggedRight\arraybackslash}p{0.390\textwidth}>{\RaggedRight\arraybackslash}p{0.140\textwidth}@{}}
\toprule\noalign{}
File & Description & Stage \\
\midrule\noalign{}
\endhead
\bottomrule\noalign{}
\endlastfoot
\href{../code/particles/particle_masses_stage5.py}{\path{particle_masses_stage5.py}} & Core spectrum prediction: gauge closure, transmutation, critical surface, Z$_6$ texture & 1--5 \\
\href{../code/particles/oph_qcd.py}{\path{oph_qcd.py}} & 4-loop MSbar $\beta$-coefficients and $\Lambda$ extraction & 6 \\
\href{../code/particles/oph_lattice_su3_quenched_v5.py}{\path{oph_lattice_su3_quenched_v5.py}} & Quenched Wilson SU(3) lattice for hadron mass ratios & 6--7 \\
\href{../code/particles/oph_predict_compare.py}{\path{oph_predict_compare.py}} & Full predictor plus PDG comparison (main entry point) & All \\
\href{../code/particles/oph_no_cheat_audit.py}{\path{oph_no_cheat_audit.py}} & Static anti-leak audit of prediction code & Audit \\
\href{../code/particles/test_oph_predict_compare.py}{\path{test_oph_predict_compare.py}} & Smoke tests plus runtime no-cheat mutation test & Tests \\
\href{../code/particles/test_particle_masses_stage5.py}{\path{test_particle_masses_stage5.py}} & Regression tests for Stage 5 predictions & Tests \\
\end{longtable}
}

\subsubsection{12.3 Running the Full Prediction}\label{123-running-the-full-prediction}

\begin{Shaded}
\begin{Highlighting}[]
\BuiltInTok{cd}\NormalTok{ code/particles/}

\CommentTok{\# Full prediction with PDG comparison table}
\ExtensionTok{python3}\NormalTok{ oph\_predict\_compare.py }\AttributeTok{{-}{-}compare}

\CommentTok{\# JSON output for programmatic use}
\ExtensionTok{python3}\NormalTok{ oph\_predict\_compare.py }\AttributeTok{{-}{-}compare} \AttributeTok{{-}{-}json}

\CommentTok{\# With internal hadron computation (slow; tiny{-}lattice demo)}
\ExtensionTok{python3}\NormalTok{ oph\_predict\_compare.py }\AttributeTok{{-}{-}compare} \AttributeTok{{-}{-}with{-}hadrons} \AttributeTok{{-}{-}hadron{-}profile}\NormalTok{ demo}
\end{Highlighting}
\end{Shaded}

\subsubsection{12.4 No-Leakage Verification}\label{124-no-leakage-verification}

\begin{Shaded}
\begin{Highlighting}[]
\BuiltInTok{cd}\NormalTok{ code/particles/}

\CommentTok{\# Static audit: checks that build\_spectrum() does not reference PDG values}
\ExtensionTok{python3}\NormalTok{ oph\_no\_cheat\_audit.py}

\CommentTok{\# Runtime mutation test: scrambles PDG table and verifies predictions unchanged}
\ExtensionTok{python3}\NormalTok{ test\_oph\_predict\_compare.py}
\end{Highlighting}
\end{Shaded}

\subsubsection{12.5 Individual Stages}\label{125-individual-stages}

\begin{Shaded}
\begin{Highlighting}[]
\BuiltInTok{cd}\NormalTok{ code/particles/}

\CommentTok{\# Stage 5 spectrum (core charged sector)}
\ExtensionTok{python3}\NormalTok{ particle\_masses\_stage5.py}

\CommentTok{\# QCD scale extraction (standalone)}
\ExtensionTok{python3}\NormalTok{ oph\_qcd.py}
\end{Highlighting}
\end{Shaded}

\subsubsection{12.5 Fetching PDG Reference Data}\label{125-fetching-pdg-reference-data}

\begin{Shaded}
\begin{Highlighting}[]
\ExtensionTok{pip}\NormalTok{ install pdg pandas}
\ExtensionTok{python3}\NormalTok{ ../tools/fetch\_pdg\_data.py}
\CommentTok{\# Output: ../pdg\_data/particle\_masses.\{csv,json\}}
\end{Highlighting}
\end{Shaded}

\begin{center}\rule{0.5\linewidth}{0.5pt}\end{center}

\subsection{Summary}\label{summary}

This part is a quantitative implementation dossier, not a second statement of the recovered-core theorem package. In the master ledger, the recovered core remains D1--D5 together with D7--D9. The role of the present part is to record what happens after that structural core is in place.

\begin{itemize}
\item \textbf{Phase I recovered core inherited from Parts I--II:} conditional relativity plus the realized Standard Model structural chain. These results are upstream prerequisites for the present file, not outputs newly established here.
\item \textbf{Phase II input-dependent / implementation layer:} D6 if the screen-capacity identification is adopted; D10 for the supplement-backed gauge/electroweak calibration sector; D11 for the Higgs/top critical-surface branch that adds no new continuous fit once D10 is fixed.
\item \textbf{Phase III continuation layer:} charged-lepton/Koide fits, texture branches, neutrino estimates beyond D6, downstream hadron extractions, proton-spin bookkeeping, and related phenomenological continuations.
\end{itemize}

The practical contribution of this part is therefore narrow but useful. It records the current D10 gauge/electroweak closure of the supplement-backed implementation, the D11 critical-surface Higgs/top branch, and a set of D12 continuations that reuse the same structural backdrop without acquiring the same claim tier.

A compact schematic of the status boundary is
\[
\text{Recovered core}
\Longrightarrow
\bigl(D6\ \text{if the capacity input is adopted}\bigr)
\Longrightarrow
D10
\Longrightarrow
D11
\Longrightarrow
D12.
\]
Here \(D10\) uses the external input \(P\), \(D6\) uses the external input \(N_{\mathrm{scr}}\), \(D11\) adds the extra critical-surface condition, and \(D12\) introduces further ans\"atze or downstream numerical machinery. Failure of a Phase-III branch does not by itself falsify the recovered core; failure of D11 challenges that secondary quantitative branch; failure of D10 challenges the current supplement-backed implementation.
